\chapter{Internal Evaluation Results}
\label{ch:results}

This chapter reports the internal results in the order needed to interpret the study. Retrieval configurations are selected first on the public-validation split (RQ3). The frozen configurations are then compared on the fixed internal held-out set (RQ1). The remaining sections localize disagreement across the recorded output views (RQ2), evaluate primary-selection interventions (RQ4), and report the internal model-scale and diagnostic-confusion analyses that contribute to RQ5. All results in this chapter measure agreement with projected benchmark labels on synthetic LingxiDiag dialogues; they are not estimates of clinical diagnostic accuracy.

\section{Retrieval-Configuration Results}
\label{sec:retrieval-results}

Table~\ref{tab:retrieval-results} reports the public-validation results used for configuration selection and the corresponding held-out results. Single and HiED--DA were selected separately because their output contracts and three-label measures differ. Single uses emitted-label hit@3, whereas HiED--DA uses genuine ranked Top-3 Accuracy. The configurations marked with $\dagger$ were frozen before the held-out set was evaluated.

\begin{table}[H]
\centering
\caption{Retrieval-strategy results on the public-validation and fixed internal held-out splits. Configurations marked with $\dagger$ were selected using validation results. Held-out results are post-selection evaluations and do not revise the retained configurations.}
\label{tab:retrieval-results}
\small

\textbf{Panel A. Single}

\vspace{0.3em}

\begin{adjustbox}{max width=\textwidth}
\begin{tabular}{|l|l|c|c|c|c|c|}
\hline
Split & Retrieval & Top-1 & Emitted hit@3 & Exact Match & Macro-F1 & Weighted-F1\\
\hline
Validation & None & 0.466 & 0.580 & 0.016 & 0.174 & 0.429\\
Validation & Global Top-5$^{\dagger}$ & 0.519 & 0.725 & 0.030 & 0.209 & 0.456\\
Validation & Parent-balanced & 0.507 & 0.741 & 0.016 & 0.195 & 0.437\\
\hline
Held-out & None & 0.466 & 0.648 & 0.017 & 0.183 & 0.426\\
Held-out & Global Top-5$^{\dagger}$ & 0.517 & 0.748 & 0.029 & 0.224 & 0.447\\
Held-out & Parent-balanced & 0.508 & 0.757 & 0.011 & 0.192 & 0.433\\
\hline
\end{tabular}
\end{adjustbox}

\vspace{0.8em}

\textbf{Panel B. HiED--DA}

\vspace{0.3em}

\begin{adjustbox}{max width=\textwidth}
\begin{tabular}{|l|l|c|c|c|c|c|}
\hline
Split & Retrieval & Top-1 & Ranked Top-3 & Exact Match & Macro-F1 & Weighted-F1\\
\hline
Validation & None & 0.509 & 0.792 & 0.456 & 0.173 & 0.428\\
Validation & Global Top-5 & 0.525 & 0.797 & 0.432 & 0.178 & 0.435\\
Validation & Parent-balanced$^{\dagger}$ & 0.523 & 0.800 & 0.428 & 0.181 & 0.440\\
\hline
Held-out & None & 0.510 & 0.786 & 0.404 & 0.174 & 0.437\\
Held-out & Global Top-5 & 0.524 & 0.815 & 0.435 & 0.196 & 0.437\\
Held-out & Parent-balanced$^{\dagger}$ & 0.518 & 0.802 & 0.409 & 0.178 & 0.434\\
\hline
\end{tabular}
\end{adjustbox}

\parbox{0.96\textwidth}{\footnotesize
\textit{Note.} Panel A reports Single's emitted-label hit@3 from the available final labels, with the primary diagnosis first. It is not a genuine ranked differential and is not used for $D_3$. Panel B reports HiED--DA's genuine ranked Top-3 Accuracy from the Diagnostician ranking. Selection used only the public-validation split and compared each architecture with itself.
}
\end{table}

\paragraph{Single.}
Global Top-5 was retained for Single. On validation, it produced the highest Top-1 Accuracy (0.519), Exact Match (0.030), Macro-F1 (0.209), and Weighted-F1 (0.456). Parent-balanced retrieval produced the highest emitted-label hit@3 (0.741), but this measure describes the final emitted labels rather than a genuine ranked differential. The held-out results show the same overall pattern: global Top-5 has the highest Top-1, Exact Match, Macro-F1, and Weighted-F1, whereas parent-balanced retrieval has the highest emitted-label hit@3. Both retrieval configurations improve held-out Top-1 over no retrieval, by 5.1 and 4.2 percentage points, respectively.

\paragraph{HiED--DA.}
Parent-balanced retrieval was the initially retained HiED--DA configuration. On validation, global Top-5 is higher on Top-1 Accuracy (0.525 versus 0.523) and Exact Match (0.432 versus 0.428). Parent-balanced retrieval is higher on genuine ranked Top-3 Accuracy (0.800 versus 0.797), Macro-F1 (0.181 versus 0.178), and Weighted-F1 (0.440 versus 0.435). Global Top-5 therefore does not satisfy the all-five promotion condition defined in Chapter~\ref{ch:experimental}, and parent-balanced retrieval remains the retained configuration.

On the held-out set, global Top-5 is higher than parent-balanced retrieval on all five HiED measures. The differences are 0.6 percentage points for Top-1, 1.3 for ranked Top-3, 2.6 for Exact Match, 1.8 for Macro-F1, and 0.3 for Weighted-F1. This validation--held-out rank reversal is a post-selection sensitivity result. It is not used to replace the configuration selected on validation.

\paragraph{Answer to RQ3.}
Retrieval provides a clear Top-1 benefit for the direct Single configuration. For HiED--DA, the differences among no retrieval, global Top-5, and parent-balanced retrieval are smaller and depend on the split and metric. The results do not establish one retrieval policy as universally superior for HiED. Parent-balanced retrieval is used in the primary HiED analyses because it is the configuration retained by the predefined validation procedure, not because it has the highest held-out point estimates.

\section{Benchmark Positioning and Main Internal Performance}
\label{sec:validation-retrieval-selection}
\label{sec:matched-baselines-retrieval}

This section addresses RQ1 on the fixed 1,000-case internal held-out set. Single uses the validation-selected global Top-5 configuration, and HiED--DA uses the validation-selected parent-balanced configuration. The comparison concerns complete systems under their stated output contracts; it is not a component-controlled test of multi-agent orchestration.
